\section{深度学习在 3D 数据上的方法}

\subsection{方法一：多视图 CNN（Multi-view CNN）}

最早将深度学习应用到 3D 数据的方法非常直接：将 3D 物体从多个角度渲染成 2D 图像，然后用已有的 2D CNN（如在 ImageNet 上预训练的模型）进行分类。

\begin{figure}[H]
\centering
\includegraphics[width=0.95\textwidth]{slides-images/slide-021.jpg}
\caption{多视图方法：将 3D 物体从多角度渲染为 2D 图像，用 2D CNN 分别处理后融合\protect\footnotemark}
\end{figure}
\footnotetext{Slides 第21页。}

\begin{knowledgebox}{多视图方法的回归}
多视图方法最初因为``不够 3D native''而被研究者抛弃，但随着 2D 图像/视频基础模型（如 V3 等）变得越来越强大，这一趋势正在回归——因为 2D 模型训练所用数据量可能比 3D 数据多百万倍。
\end{knowledgebox}

\subsection{方法二：体素 CNN（Volumetric CNN）}

最直接的 3D 原生方法：将 2D 卷积扩展为 3D 卷积，在体素数据上运行。

\begin{figure}[H]
\centering
\includegraphics[width=0.95\textwidth]{slides-images/slide-022.jpg}
\caption{3D CNN / 体素卷积网络：将 2D 卷积核扩展为 3D，直接在体素数据上运行\protect\footnotemark}
\end{figure}
\footnotetext{Slides 第22页。}

\begin{warningbox}{体素方法的计算瓶颈}
体素分辨率为 $N^3$，将 2D 的 $N^2$ 复杂度提升了一个维度。当 $N=256$ 时，体素需要 $256^3 \approx 1700$ 万个元素，内存和计算量巨大。而且大量体素对应空白区域，造成严重浪费。为缓解这一问题，人们开发了\textbf{稀疏卷积}（Sparse Convolution）和\textbf{八叉树}（Octree）等技术。
\end{warningbox}

\begin{figure}[H]
\centering
\includegraphics[width=0.95\textwidth]{slides-images/slide-024.jpg}
\caption{八叉树（Octree）：自适应分辨率的层级结构，仅在有物体的区域细分，大幅减少计算量\protect\footnotemark}
\end{figure}
\footnotetext{Slides 第24页。}

\subsection{方法三：PointNet --- 直接处理点云}

\textbf{PointNet}（Stanford, Qi et al. 2017）是 3D 深度学习领域的里程碑工作，提出了直接在点云上运行的神经网络。

\begin{figure}[H]
\centering
\includegraphics[width=0.95\textwidth]{slides-images/slide-027.jpg}
\caption{PointNet 架构：对每个点独立计算特征嵌入，然后用对称函数（max/sum）聚合为全局特征\protect\footnotemark}
\end{figure}
\footnotetext{Slides 第27页。}

PointNet 需要解决两个关键问题：

\begin{importantbox}{PointNet 的两大设计要求}
\begin{enumerate}
    \item \textbf{排列不变性}（Permutation Invariance）：点的编号顺序不影响输出
    \item \textbf{采样不变性}（Sampling Invariance）：不同采样密度下输出应一致
\end{enumerate}
解决方案极其简洁：对每个点独立计算嵌入，然后用\textbf{对称函数}（如 max pooling 或 sum）聚合所有点的特征。
\end{importantbox}

PointNet 的工作流程：
\begin{enumerate}
    \item 输入 $N$ 个点，每个点有 $(x, y, z)$ 坐标
    \item 对每个点独立地通过 MLP 计算高维嵌入
    \item 用 max pooling 聚合所有点的特征为一个全局向量
    \item 通过全连接层输出分类结果
\end{enumerate}

后续改进包括 PointNet++（引入局部特征）、图神经网络（将点作为图节点、邻近关系作为边）等。

\subsection{点云的距离度量}

对于点云生成任务，需要定义两个点云之间的``距离''。两种常用度量：

\begin{figure}[H]
\centering
\includegraphics[width=0.95\textwidth]{slides-images/slide-029.jpg}
\caption{Chamfer Distance 与 Earth Mover's Distance：两种点云间距离度量\protect\footnotemark}
\end{figure}
\footnotetext{Slides 第29页。}

\textbf{Chamfer Distance}：
$$d_{CD}(S_1, S_2) = \sum_{x \in S_1} \min_{y \in S_2} \|x - y\|^2 + \sum_{y \in S_2} \min_{x \in S_1} \|x - y\|^2$$

\textbf{Earth Mover's Distance (EMD)}：
$$d_{EMD}(S_1, S_2) = \min_{\phi: S_1 \to S_2} \sum_{x \in S_1} \|x - \phi(x)\|$$

\begin{knowledgebox}{两种距离的区别}
Chamfer Distance 计算简单但对异常值敏感；EMD 需要求解最优传输问题，计算复杂但更鲁棒。实际应用中，Chamfer Distance 因计算效率更常被使用。
\end{knowledgebox}

\subsection{本章小结}

\begin{itemize}
    \item 多视图 CNN：利用强大的 2D 预训练模型，但丢失了 3D 原生信息
    \item 体素 CNN：最直接的 3D 扩展，但面临 $O(N^3)$ 的计算瓶颈
    \item PointNet：通过排列不变的对称函数，开创了直接处理点云的范式
    \item 不同方法各有优劣，选择取决于数据格式和任务需求
\end{itemize}

%% ========== 神经隐式表示 ==========

